\section{Introduction}
\label{sec:introduction}

The practical policy concern behind this thesis is not simply that pricing algorithms can raise prices in a laboratory model. It is whether the messy information and timing conditions of real markets make that concern smaller or larger. Algorithmic pricing is already discussed by competition authorities in sectors such as travel, entertainment, retail, and fuel, and empirical work studies its effects in online retail and German petrol markets \citep{oecd2025,brown2023,assad2024}. In the German petrol market, for example, \citet{assad2024} find that markups rose around the time major chains adopted algorithmic pricing software. These markets do not look like frictionless textbook games. Firms observe rivals through booking systems, web scrapers, price feeds, shelf tags, and reporting platforms, all of which can be delayed, noisy, or tied to fixed update schedules. If those frictions already weaken algorithmic collusion, then the policy problem is less urgent than the clean simulation literature suggests. If they can instead preserve or even restore collusion when combined, then the problem is larger than one would infer from studying each friction in isolation.

The legal difficulty is that cartel enforcement is built around agreement. Explicit price fixing is illegal because it raises prices above the competitive level, transfers surplus from consumers to firms, and creates deadweight loss \citep{businessgovnl}. But autonomous algorithmic coordination may not involve the kind of communication, meeting of minds, or explicit agreement that competition law is normally designed to prove. The OECD therefore treats algorithmic collusion as a challenge for traditional concepts of agreement and tacit coordination \citep{oecd2017,oecd2025}. \citet{calvano2020} make the concern concrete: two independent Q-learning algorithms in a repeated pricing game learn prices about 84\% of the way from competitive Nash to joint monopoly, without ever communicating. In economic terms the outcome resembles a cartel, but the conduct is generated by trial-and-error learning rather than an explicit agreement. The full literature, including the follow-up simulation studies and the empirical evidence, is reviewed in Section~\ref{sec:literature}.

The natural next question is whether the Calvano result survives once the information environment becomes less clean. This matters both descriptively and for policy. Descriptively, real pricing algorithms do not always see rival prices instantly or perfectly. They often receive stale, rounded, scraped, or otherwise imperfect signals. For policy, these same features are possible indirect levers: a regulator may not be able to prosecute communication-free learning directly, but could still ask whether delayed disclosure, noisy public price feeds, or commitment windows would make algorithmic coordination less likely. Existing papers study some of these channels one at a time. What is less clear is whether they still work when imposed together, or whether one friction can undo the effect of another.

I keep the \citet{calvano2020} pricing environment fixed and change only the way the algorithms observe or time their decisions. In the baseline, two firms repeatedly choose prices, consumers respond according to logit demand, and each firm is controlled by a Q-learning algorithm that learns from its own profits and the rival's previous price. I measure the outcome using the profit-gain index $\Delta$. This is an economic collusion score: $\Delta=0$ means profits are at the competitive Nash benchmark, while $\Delta=1$ means profits are at the joint-monopoly benchmark. A value such as $\Delta=0.5$ therefore means that the algorithms have recovered about half of the profit gap between competition and joint monopoly, which is already a substantial supracompetitive outcome. The thesis asks how this score changes when three frictions are introduced separately and jointly:
\begin{enumerate}[label=\textbf{RQ\arabic*.}, leftmargin=3em, labelsep=0.6em, itemsep=5pt, parsep=0pt, topsep=10pt, partopsep=3pt]
    \item How does observation latency (delayed observation of the rival's price) affect collusion?
    \item How does noisy monitoring (Gaussian misobservation of the rival's price) affect collusion?
    \item How does asynchronous price setting (Maskin--Tirole alternation between movers) affect collusion?
    \item Do these frictions interact non-linearly when imposed jointly?
\end{enumerate}

Varying each friction separately gives three different response patterns in the collusion score: latency produces a sharp initial drop, noise produces a smooth decline toward the competitive benchmark, and asynchrony produces a one-time drop followed by a plateau. A $2\times 2\times 2$ factorial experiment, replicated at three session counts to confirm robustness, then identifies two interaction patterns. Latency and noise compound: imposing both jointly destroys more collusion than independent dampening would predict, because the two frictions degrade different dimensions of the rival's observation, timing and value. Asynchrony has the opposite interaction. Although it reduces collusion on its own, it partially restores collusion when combined with latency or noise. The policy relevance is therefore cautious rather than prescriptive. The results do not say that a particular regulatory package should be adopted. They show that the case for intervention cannot be assessed by testing one friction at a time, because realistic combinations can amplify some anti-collusion effects while offsetting others.
